% ============================================================
% APPENDIX — 9  sections
% A: Statistical Significance
% B: Evaluated Models and Tasks
% C: Ablation Studies
% D: Per-Category Rankings and Figures
% E: Qualitative Error Analysis
% F: Extended Related Work
% G: Extended Limitations
% H: Cost and Throughput Details
% I: Reproducibility
% ============================================================


% ============================================================
\section{Statistical Significance Testing}
\label{app:significance}
% ============================================================

\paragraph{Bootstrap procedure.}
We implement a paired bootstrap test \citep{efron1979bootstrap} comparing Gemini~3.1~Pro against the best embedding model in each comparison.
The resampling unit is the task level: each bootstrap replicate draws \mteblm{} tasks with replacement from the full \mteblm{} task set, and the score difference is computed as the mean-of-means difference across the resampled task set.
Task-level resampling accommodates the different metrics used across categories (accuracy, Spearman $\rho$, V-measure, average precision, and Recall@1). Equal task weighting matches the macro-average score.

We report unadjusted $p$-values for five planned category-level comparisons. Under a Bonferroni threshold of $\alpha$\,=\,0.01, classification remains significant and retrieval falls above the threshold.

Table~\ref{tab:significance} reports the full results.
The overall comparison (Pro vs.\ Octen-8B across \mteblm) yields $\Delta$\,=\,$+$0.3, 95\%~CI\,=\,[$-$2.4, +3.1], $p$\,=\,0.85, a statistical tie.
Category-level tests identify two significant differences:
\begin{itemize}
    \item \textbf{Classification:} embeddings win ($\Delta$\,=\,$-$5.6, CI\,=\,[$-$9.2, $-$2.4], $p$\,$<$\,0.01)
    \item \textbf{Retrieval:} LLMs win ($\Delta$\,=\,+8.5, CI\,=\,[+0.2, +16.8], $p$\,$<$\,0.05)
\end{itemize}
STS ($p$\,=\,0.74), clustering ($p$\,=\,0.97), and pair classification ($p$\,=\,0.49) show no significant paradigm difference.

\input{tables/significance}


% ============================================================
\section{Evaluated Models and Tasks}
\label{app:models_detail}
% ============================================================

This section lists the evaluated models, category- and task-level results, and benchmark tasks and metrics.

\subsection{Model Overview}

Table~\ref{tab:full_models} lists all 36 evaluated models — 10 LLMs and 26 embedding models — with parameter counts, mean \mteblm{} scores, total benchmark costs, and paper references.
Models range from mE5-small (118M parameters, \$0.001) to Gemini~3.1~Pro (\$154.14).
Models without a dedicated paper are marked ``--''.

\input{tables/full_models}

\subsection{Full Per-Category Rankings}
\label{app:models}

Table~\ref{tab:category_scores} compares Pro with the best embedding in each category. Table~\ref{tab:full_results_by_category} gives per-category scores for all 36 models.
Pro ranks first overall largely because of retrieval; on classification, it ranks below ten embedding models.
The reasoning-capable Gemini models rank far above the non-reasoning Flash-Lite, which places 31st of 36 overall.

\input{tables/category_scores}
\input{tables/full_results_by_category}

\subsection{Per-Task Scores (Representative Models)}
\label{app:per_task}

Table~\ref{tab:per_task_scores} reports individual task scores for 5 representative models (2 LLMs, 3 embeddings) across all \mteblm{} tasks, alongside the best embedding score per task.
The scores vary substantially by task. For example, Pro is strongest on FQuAD retrieval, and SFR-2 is strongest on SprintDuplicateQuestions.
This complementarity supports the hybrid-pipeline recommendation in \S\ref{sec:discussion}.

\input{tables/per_task_scores}

\subsection{Full Per-Task Scores for All Models}
\label{app:full_per_task_all}

Tables~\ref{tab:full_per_task_all_models}--\ref{tab:full_per_task_ret} provide the complete score matrix for all 36 models across every \mteblm{} task, presented one table per category so that each fits the page upright.
The tables share a common layout: models are rows, ordered by overall \mteblm{} score, with the ten LLMs listed first and the 26 embedding models below; the category's tasks are columns, with its metric given in the caption.
The final column gives each model's mean over the category, matching the corresponding column of Table~\ref{tab:main_results}.
Bold marks the best value in a column, whether a single task or the category mean.
Because rows are ordered by overall score rather than by category, this column is not monotonic: it shows where a model over- or under-performs its overall rank.
These tables extend Table~\ref{tab:per_task_scores} to all models.

\input{tables/full_per_task_all_models}

\subsection{Benchmark Task Suite}
\label{app:datasets}

Table~\ref{tab:datasets} lists all 37 tasks with languages, sample counts, and source citations.
All are held-out subsets (seed\,=\,42) derived from MTEB and MMTEB tasks, hosted on Hugging Face under \texttt{mteb/llm-eval-*}; the exact dataset path and pinned revision for each task are listed in the released repository.
Table~\ref{tab:token_budget} provides per-task token budgets for cost estimation, broken down by split.

\input{tables/datasets}

\input{tables/token_budget}


% ============================================================
\section{Ablation Studies}
\label{app:ablation_detail}
% ============================================================

This section details the reduced-thinking and few-shot experiments and provides a per-task retrieval breakdown.

\subsection{Reduced-Thinking Ablation Details}
\label{app:ablation_reduced}

Our reduced-thinking ablation uses Gemini~3~Flash with \texttt{reasoning\_effort=low}, which instructs the model to minimise internal chain-of-thought reasoning.
We evaluate on all 6 retrieval tasks and 3 representative classification tasks (IMDB, Banking77, ToxicConversations).

The reduction differs by task type: 54--94\% on retrieval and 12--30\% on classification.
This suggests the model adaptively allocates more reasoning to retrieval (where documents must be read and compared) than classification (where the decision is more immediate).

Reduced reasoning preserves or improves all six Flash retrieval scores and four of six cross-family averages (Figure~\ref{fig:cost_breakdown}b).
We hypothesise that excessive reasoning introduces second-guessing: the model considers multiple interpretations of a document's relevance when a more direct reading would suffice.

On classification, all three tasks show $|\Delta| < 1.0$, indicating no measurable benefit from additional reasoning in this ablation.
This result is consistent with the supervision asymmetry between kNN and zero-shot classification.

\input{tables/ablation_nocot}

\input{tables/ablation_fewshot}

\subsection{Per-Task Retrieval Analysis}
\label{app:retrieval_detail}

The six retrieval tasks in our benchmark span qualitatively different retrieval challenges; per-task scores for all 36 models are in Table~\ref{tab:full_per_task_ret}.
We provide a per-task breakdown of the aggregate retrieval advantage ($+$8.5) reported in \S\ref{sec:categories}.
Pro leads five retrieval tasks, and Octen-8B leads AILAStatutes.

\paragraph{FQuAD (French reading-comprehension QA).}
Answering requires reading French passages in context and matching the supporting one.
Pro (92.0) leads the best embedding (72.0) by a wide margin, the largest LLM advantage in the suite.

\paragraph{Consumer-Contracts QA (legal).}
Yes/no questions over consumer-contract clauses.
Pro (86.0) outperforms the best embedding (70.0), consistent with the need to interpret the clause.

\paragraph{Public-Health QA.}
Pro (90.0) edges the best embedding (85.0); factual health QA is partly captured by dense retrieval, narrowing the gap.

\paragraph{HC3-Finance QA.}
Pro (71.0) narrowly leads the best embedding (67.0) on financial QA retrieval.

\paragraph{TwitterHjerne (Danish social media).}
Retrieving relevant tweets given a query.
Pro (33.4) slightly outperforms embeddings (best: 31.5); both handle keyword-centric retrieval similarly.

\paragraph{AILAStatutes (legal statute retrieval): where reasoning hurts.}
Matching case facts to relevant statutes.
The best embedding (Octen-8B, 23.2) outperforms Pro (14.5): Pro over-retrieves broadly associated provisions, while cosine similarity naturally produces a narrower, more precise ranking.



% ============================================================
\section{Per-Category Rankings and Figures}
\label{app:category_rankings}
% ============================================================

This section provides task- and model-level figures, followed by category-level rankings.
Figure~\ref{fig:taskrange} is the task-level view of the paper's central claim and is referenced from \S\ref{sec:categories}.

\subsection{Task- and Model-Level Figures}
\label{app:figures}

The figures below show task-level performance, throughput, category profiles, and cross-model correlations.

\begin{figure*}[htbp]
\centering
\includegraphics[width=\textwidth]{task_performance_range.pdf}
\caption{\textbf{Per-task performance.}
Horizontal bars show the full embedding score range (min to max across 26 models) for each \mteblm{} task, grouped by category.
Circle = best embedding; diamond = best LLM on that task, taken over all ten (six different LLMs hold it across the suite, Gemini~3.1~Pro on 21 of 37 tasks).
Where the diamond falls inside the bar, some embedding model already matches the best LLM: this holds on 7 of 8 classification tasks and 7 of 10 STS tasks, but on only 1 of 6 retrieval tasks.}
\label{fig:taskrange}
\end{figure*}

\begin{figure}[htbp]
\centering
\includegraphics[width=\linewidth]{throughput_comparison.pdf}
\caption{\textbf{Same-hardware inference throughput.}
Two open-weight LLMs and seven representative embedding models (118M--14B) served on one H100 (tokens/second, log scale); the full 26-model embedding throughput is in Table~\ref{tab:embedding_throughput}.
Even the slowest embedding runs $\sim$2.5$\times$ faster than the fastest LLM; the fastest runs $\sim$736$\times$ faster.}
\label{fig:throughput}
\end{figure}

\begin{figure}[htbp]
\centering
\includegraphics[width=0.85\linewidth]{radar_capability_profile.pdf}
\caption{\textbf{Capability profiles across five task categories.}
Gemini~3.1~Pro (red), Qwen3.6-27B (amber), and Octen-8B (blue).
The LLMs lead on retrieval, while Octen-8B leads on classification and remains competitive across categories.}
\label{fig:radar}
\end{figure}

\begin{figure}[htbp]
\centering
\includegraphics[width=\linewidth]{task_correlation_heatmap.pdf}
\caption{\textbf{Cross-model task-behaviour correlation.}
Pearson correlations based on per-task \mteblm{} scores.
LLMs and embedding models form visually distinct clusters: LLMs are strongly intercorrelated while embedding models cluster by architecture family.
Lower inter-paradigm correlations are consistent with their different category profiles.}
\label{fig:heatmap}
\end{figure}

\FloatBarrier

\subsection{Per-Category Model Rankings}
\label{app:rankings_figs}

Figures~\ref{fig:ranking_overall}--\ref{fig:ranking_ret} rank all 36 evaluated models overall and within each of the five MTEB task categories.
They show category-dependent performance: LLMs lead on retrieval, embedding models lead on classification, and the two paradigms are statistically indistinguishable on clustering, STS, and pair classification.
Flash-Lite ranks near the bottom in every category, below the reasoning-capable Gemini models.

\begin{figure*}[htbp]
\centering
\includegraphics[width=\textwidth]{ranking_Overall.pdf}
\caption{\textbf{Overall rankings.} Pro (77.6) narrowly leads Octen-8B (77.2); the difference lies within statistical noise ($p$\,=\,0.85). Flash-Lite (64.5) ranks near the bottom.}
\label{fig:ranking_overall}
\end{figure*}

\begin{figure*}[htbp]
\centering
\includegraphics[width=\textwidth]{ranking_Classification.pdf}
\caption{\textbf{Classification rankings.} SFR-2 ranks first (90.8); Pro scores 85.2 and ranks below ten embedding models.}
\label{fig:ranking_clf}
\end{figure*}

\begin{figure*}[htbp]
\centering
\includegraphics[width=\textwidth]{ranking_Clustering.pdf}
\caption{\textbf{Clustering rankings.} Statistical tie; SFR-2 (66.7) edges Pro (66.6). Flash-Lite scores 21.7.}
\label{fig:ranking_clu}
\end{figure*}

\begin{figure*}[htbp]
\centering
\includegraphics[width=\textwidth]{ranking_STS.pdf}
\caption{\textbf{STS rankings.} Qwen3-E-4B (88.8) and Pro (88.5) are essentially tied; no significant paradigm difference ($p$\,=\,0.74).}
\label{fig:ranking_sts}
\end{figure*}

\begin{figure*}[htbp]
\centering
\includegraphics[width=\textwidth]{ranking_PairClassification.pdf}
\caption{\textbf{Pair Classification rankings.} KaLM-12B (87.1) outperforms Pro (83.2).}
\label{fig:ranking_pair}
\end{figure*}

\begin{figure*}[htbp]
\centering
\includegraphics[width=\textwidth]{ranking_Retrieval.pdf}
\caption{\textbf{Retrieval rankings.} Pro leads the best embedding (64.5 vs.\ 56.0) and ranks first on five of six tasks.}
\label{fig:ranking_ret}
\end{figure*}

\FloatBarrier

\subsection{Full Retrieve-then-Rerank Matrix}
\label{app:rerank_matrix}

Table~\ref{tab:reranker_matrix} reports the complete matrix summarised by Figure~\ref{fig:rerank} (\S\ref{sec:rerank}): all four first-stage retrievers crossed with every cross-encoder and LLM listwise reranker, on BEIR and BRIGHT.

\input{tables/reranker_matrix}


% ============================================================
\section{Qualitative Error Analysis}
\label{app:samples}
% ============================================================

We present model predictions from the evaluation files to illustrate differences between LLM and embedding approaches.
The examples cover classification ambiguity, STS calibration, retrieval precision, and differences between Pro and Flash.

\subsection{Classification: Over-Interpretation vs.\ Anchored Reference}

Several classification errors involve ambiguous inputs for which the LLM infers an intent that differs from the dataset convention.
With kNN, an embedding classifier can anchor its prediction to similar labelled examples, which helps on fine-grained tasks such as Banking77 and MassiveIntent.

\paragraph{Banking77 (77 classes): ``locate my card.''}
\begin{quote}
\small
\textbf{Input:} ``How do I locate my card?''

\textbf{Pro reasoning:} ``The user is asking how to locate their card, indicating that they have misplaced it or it is missing. This aligns with the intent for a lost or stolen card.''

\textbf{Pro prediction:} \texttt{lost\_or\_stolen\_card} \hfill \textcolor{red}{Incorrect (gold: \texttt{card\_arrival})}

\textbf{Flash prediction:} \texttt{card\_arrival} \hfill \textcolor{green!60!black}{Correct}
\end{quote}

\noindent Flash predicts the gold label; Pro predicts \texttt{lost\_or\_stolen\_card}.
The LLM infers ``locate'' implies loss; Flash's lower-depth reasoning follows the surface-level reading.
In the training data, ``locate my card'' aligns with delivery tracking.
The kNN classifier matches to examples like ``Where is my card?'' and correctly predicts \texttt{card\_arrival}.

\paragraph{Banking77: ``know when my card will arrive.''}
\begin{quote}
\small
\textbf{Input:} ``Is there a way to know when my card will arrive?''

\textbf{Pro reasoning:} ``The user is asking for information regarding the timeframe or expected date for their card's delivery, which aligns with the \texttt{card\_delivery\_estimate} category.''

\textbf{Pro prediction:} \texttt{card\_delivery\_estimate} \hfill \textcolor{red}{Incorrect (gold: \texttt{card\_arrival})}

\textbf{Flash prediction:} \texttt{card\_delivery\_estimate} \hfill \textcolor{red}{Incorrect (gold: \texttt{card\_arrival})}
\end{quote}

\noindent Both models make the same error here, illustrating the ambiguity between ``arrival'' and ``delivery estimate'' without labelled context.
The kNN classifier uses the full 77-class reference set to resolve such near-neighbour cases.

\paragraph{MassiveIntent (60 classes): ``remind me about my schedule.''}
\begin{quote}
\small
\textbf{Input (en):} ``remind me about my schedule for the afternoon''

\textbf{Pro reasoning:} ``The user is asking to be reminded about their schedule for the afternoon, which involves checking or querying their calendar events.''

\textbf{Pro prediction:} \texttt{calendar\_query} \hfill \textcolor{red}{Incorrect (gold: \texttt{calendar\_set})}
\end{quote}

\noindent ``Remind me'' is an action verb (set a reminder) but the LLM interprets it as information-seeking (query).
This example illustrates the ambiguity between action and query intents.

\paragraph{MassiveScenario (German): ``wir brauchen rosa.''}
\begin{quote}
\small
\textbf{Input (de):} ``wir brauchen rosa'' (we need pink)

\textbf{Pro reasoning:} ``The utterance 'wir brauchen rosa' (we need pink) implies adding an item ('rosa' / pink) to a shopping list or note-taking application.''

\textbf{Pro prediction:} \texttt{lists} \hfill \textcolor{red}{Incorrect (gold: \texttt{iot})}
\end{quote}

\noindent The LLM correctly reasons about the surface form but miscategorises the scenario: \texttt{iot} refers to controlling physical smart-home devices, not arbitrary colour references.

\begin{quote}
\small
\textbf{Input (de):} ``wie ist die einstellung f{\"u}r die helligkeit an meinem bildschirm'' (what is the brightness setting on my screen)

\textbf{Pro reasoning:} ``The user is asking about the brightness setting on their screen, which involves controlling or checking the state of a device. This falls under the IoT category.''

\textbf{Pro prediction:} \texttt{iot} \hfill \textcolor{red}{Incorrect (gold: \texttt{general})}
\end{quote}

\noindent Both errors reflect the fine-grained scenario taxonomy: \texttt{iot} refers to controlling physical smart-home devices, not arbitrary brightness settings or colour references.
The LLM correctly identifies device-related surface features but miscategorises the scenario.

\paragraph{AmazonCounterfactual: hypothetical ``would be great.''}
\begin{quote}
\small
\textbf{Input:} ``Would be great to have on my tablet instead of carrying the book.''

\textbf{Pro reasoning:} ``The review expresses a hypothetical situation ('have on my tablet') that has not actually happened, indicated by the phrase 'Would be great'. This describes an alternative condition and its potential consequence.''

\textbf{Pro prediction:} \texttt{counterfactual} \hfill \textcolor{red}{Incorrect (gold: \texttt{not-counterfactual})}
\end{quote}

\noindent The phrase ``would be'' supports the LLM's hypothetical reading. The dataset convention labels forward-looking wishes as \texttt{not-counterfactual}.
This may reflect a mismatch between a plausible linguistic interpretation and the dataset's annotation convention.

\paragraph{When LLMs succeed on classification.}
\begin{quote}
\small
\textbf{Input:} ``I still have not received my new card, I ordered over a week ago.''

\textbf{Pro prediction:} \texttt{card\_arrival} \hfill \textcolor{green!60!black}{Correct}
\end{quote}

\noindent In this less ambiguous example, the LLM matches the gold annotation.


\subsection{STS: Calibration and Semantic Precision}

The examples below show two kinds of disagreement between LLM scores and STS annotations: sensitivity to fine distinctions and higher scores for inferred mechanistic equivalence.

\paragraph{STSBenchmark: agreement on a simple paraphrase.}
\begin{quote}
\small
\textbf{S1:} ``open again, open anew'' \quad \textbf{S2:} ``open again or anew.''

\textbf{Pro:} 5.0 \quad \textbf{Gold:} 5.0 \quad \textbf{$|\Delta|$:} 0.0
\end{quote}

\noindent The LLM score matches the annotation for this near-identical paraphrase.

\paragraph{STSBenchmark: conservative scoring on ideological distinctions.}
\begin{quote}
\small
\textbf{S1:} ``\ldots a shift from a muslim identity to an islamist identity.''

\textbf{S2:} ``\ldots to change a muslim identity for an islamic identity.''

\textbf{Pro:} 4.5 \quad \textbf{Gold:} 5.0 \quad \textbf{$|\Delta|$:} 0.5
\end{quote}

\noindent The LLM distinguishes ``Islamist'' (political ideology) from ``Islamic'' (religious adjective); the gold standard treats them as equivalent.
This linguistically defensible distinction reduces agreement with the benchmark annotation.

\paragraph{BIOSSES: conservative scoring on biomedical paraphrases.}
\begin{quote}
\small
\textbf{S1:} ``It has recently been shown that Craf is essential for Kras G12D-induced NSCLC.''

\textbf{S2:} ``It has recently become evident that Craf is essential for the onset of Kras-driven non-small cell lung cancer.''

\textbf{Pro reasoning:} ``Both sentences state that Craf is essential for the development of Kras-driven NSCLC. Sentence 1 specifies the particular mutation (Kras G12D), whereas Sentence 2 is slightly more general.''

\textbf{Pro:} 3.8 \quad \textbf{Gold:} 4.0 \quad \textbf{$|\Delta|$:} 0.2
\end{quote}

\noindent The LLM correctly notes the specificity difference (``Kras G12D'' vs.\ ``Kras-driven'') and downweights accordingly.
The gold annotation gives 4.0, treating the two as highly similar; the LLM's 3.8 is slightly more conservative.

\paragraph{BIOSSES: LLM over-scores on mechanistic equivalence.}
\begin{quote}
\small
\textbf{S1:} ``More recently, IDH mutations and resultant 2-hydroxyglutarate (2HG) production in leukemia cells were found to impair TET2 function.''

\textbf{S2:} ``It has also been recently reported that mutations of the isocitrate dehydrogenase genes IDH1 and IDH2 inhibit TET catalytic activity.''

\textbf{Pro reasoning:} ``Both sentences describe the exact same biochemical pathway: mutations in IDH genes lead to the production of 2-hydroxyglutarate (2-HG), which in turn inhibits TET2 enzymatic activity.''

\textbf{Pro:} 3.8 \quad \textbf{Gold:} 3.2 \quad \textbf{$|\Delta|$:} 0.6
\end{quote}

\noindent The LLM identifies a mechanistic relationship and scores the pair 3.8, compared with the gold score of 3.2.
Expert adjudication would be needed to determine whether this difference reflects model calibration or annotation criteria.


\subsection{Retrieval: Reasoning Depth vs.\ Geometric Precision}

The retrieval examples show how joint query--document processing can help on reading-comprehension tasks and how broader associations can hurt precise statute retrieval.

\paragraph{FQuAD: reading comprehension in context.}
On French reading-comprehension retrieval, the LLM reads each candidate passage in the prompt and selects the one that answers the query, reaching Recall@1\,=\,92.0 versus 72.0 for the best embedding.
The LLM processes the query jointly with the candidate passages, which may explain part of the observed advantage.

\paragraph{AILAStatutes: where reasoning hurts.}
\begin{quote}
\small
\textbf{Query:} Legal case about a bank employee terminated for misconduct.

\textbf{Pro:} Returns 14 documents spanning broad legal concepts (contract law, employment law, due process). Gold: only 4 specific statutes.
The LLM's associative reasoning retrieves plausible but non-relevant provisions.

\textbf{Octen-8B:} Returns a narrower, more precise set via cosine similarity with legal embeddings.
\end{quote}

\noindent In this example, Pro retrieves plausible but irrelevant provisions, while the embedding model produces a narrower ranking.

\paragraph{TwitterHjerne: corpus-in-context precision.}
\begin{quote}
\small
\textbf{Query (Danish):} ``S{\o}nnen vil gerne lave \#pebern{\o}dder. De par gange jeg har pr{\o}vet det, blev de kun OK. Er der nogen, der kan anbefale en opskrift? \#twitterhjerne''
(``My son would like to make peppernuts. The few times I've tried, they were only OK. Can anyone recommend a recipe? \#twitterhjerne'')

\textbf{Pro:} Returns documents [1, 2]: exact match. \textbf{Recall@1\,=\,1.0 on this query.}

\textbf{Best embedding:} Recall@1\,=\,31.5 over the full task.
\end{quote}

\noindent For this query with explicit topical keywords, corpus-in-context retrieval returns the correct tweet first.
The overall task-level advantage for Pro (33.4 vs.\ 31.5 for embeddings) is modest, reflecting that both paradigms handle keyword-centric retrieval similarly.


\subsection{Model Scale Effects: Flash vs.\ Pro on the Same Input}

\paragraph{Banking77: Flash succeeds where Pro over-interprets.}

As shown above, Flash correctly predicts \texttt{card\_arrival} for ``How do I locate my card?'' while Pro predicts \texttt{lost\_or\_stolen\_card}.
In this example, more reasoning does not improve the prediction.
Pro's extended chain-of-thought introduces an interpretive step (``locate implies misplacement'') that Flash's lighter-weight reasoning bypasses.
Aggregate scores provide context: Pro remains slightly stronger than Flash on classification overall.

\paragraph{Retrieval scale effects.}

On FQuAD reading-comprehension retrieval, all three Gemini models outscore the best embedding, with a clear quality ladder:
\begin{itemize}
    \item Pro (Recall@1\,=\,92.0): reads and matches passages precisely.
    \item Flash (Recall@1\,=\,88.0): captures most of the advantage.
    \item Flash-Lite (Recall@1\,=\,84.0): still well above embeddings, even without reasoning.
    \item Best embedding (Recall@1\,=\,72.0): limited to single-vector matching.
\end{itemize}

Flash-Lite already outperforms the best embedding on this task; Pro adds a modest gain at much higher cost. This result is consistent with the reduced-thinking ablation (\S\ref{sec:ablation}).

\paragraph{Summary.}
The examples illustrate four patterns:
\begin{enumerate}
\item \textbf{Ambiguous label boundaries}: LLM predictions can differ from dataset conventions; kNN anchors decisions to labelled examples.
\item \textbf{STS calibration}: model scores can differ from annotations when they attend to different semantic details.
\item \textbf{Retrieval breadth}: joint processing can help reading comprehension but can also retrieve broadly related material.
\item \textbf{Model scale}: the value of additional capacity or reasoning varies across examples and must be weighed against cost.
\end{enumerate}
Estimating the frequency of these patterns requires a systematic error analysis.


% ============================================================
\section{Extended Related Work}
\label{app:extended_related}
% ============================================================

\subsection{LLMs as Text Encoders and Judges}

The boundary between generative LLMs and discriminative embedding models has blurred considerably.
LLM-as-a-judge frameworks \citep{zheng2023judging} repurpose frozen LLMs for evaluation. \citet{sun2023textclassificationlargelanguage} find that zero-shot LLM classification underperforms fine-tuned models on standard benchmarks and propose chain-of-thought prompting to narrow the gap.
Fine-tuned encoder models outperform zero-shot frontier LLMs across the classification benchmarks tested by \citet{bucher2024finetuned}, with margins of 5 to over 75 F1 points depending on task granularity. This is consistent with our 5.6-point advantage for kNN-equipped embedding pipelines.
The LOFT benchmark \citep{lee2024longcontextlanguagemodelssubsume} evaluates long-context LLMs on retrieval by placing entire corpora in the prompt (corpus-in-context prompting); our retrieval evaluation adopts the same protocol.

On the representational side, several lines of work extract embeddings from LLMs.
SGPT \citep{muennighoff2022sgptgptsentenceembeddings} uses weighted mean pooling over GPT decoder states. \citet{jiang2023scalingsentenceembeddingslarge} use in-context learning to extract sentence embeddings from decoder-only LLMs, while \citet{behnamghader2024llm2veclargelanguagemodels} adapt decoder-only LLMs into sentence encoders. \citet{springer2025repetitionimproveslanguagemodel} show that repeating the input can improve these embeddings by enabling effective bidirectional attention.
GritLM \citep{muennighoff2025generativerepresentationalinstructiontuning} goes further, unifying generation and representation in a single model that achieved strong performance on both MTEB and generative benchmarks at the time of its release.

These methods combine elements of both paradigms. At inference time, they produce fixed-length vectors compared via cosine similarity, so our cost accounting treats them as embedding models.
Several of our 26 evaluated models (e.g., E5-Mistral, GritLM) are built on LLM backbones and incorporate LLM-derived representations.

HUME \citep{assadi2025humemeasuringhumanmodelperformance} also benchmarks LLMs as annotators on MTEB tasks and finds that they fall short of human judgment on reranking. This complements our finding that zero-shot LLM classifiers underperform embedding pipelines with kNN access to labelled references.
Our study adds a controlled comparison on identical tasks, data, and metrics, together with exact cost accounting.

\subsection{Frontier LLMs and Open-Weight Reasoning Models}

LLMs have advanced rapidly since GPT-4 \citep{openai2024gpt4technicalreport}.
OpenAI's o1 \citep{openai2024o1} demonstrated extended chain-of-thought reasoning at inference time using reinforcement learning.
At the time of our evaluation, closed-source models with controllable reasoning included Gemini~3.1~Pro, GPT-5 \citep{singh2026openaigpt5card}, and Claude~Opus~4.6 \citep{anthropic2026claudeopus46}. These models support extended thinking tokens billed separately from output; newer models have since appeared (e.g.\ GPT-5.4).
Open-weight alternatives include Llama~3 \citep{grattafiori2024llama3herdmodels}, Qwen3 \citep{yang2025qwen3technicalreport}, and DeepSeek-R1 \citep{deepseekai2025deepseekr1incentivizingreasoningcapability}. DeepSeek-R1 shows that reinforcement-learning-based training can elicit strong reasoning at lower inference cost.
Gemma~3 \citep{gemmateam2025gemma3technicalreport} is a compact multilingual model family.

We evaluate ten LLMs across six families (\S\ref{sec:setup}), including DeepSeek-R1 and open-weight Qwen3.6, GLM, Kimi, and MiniMax models alongside Gemini~3. The cost and task-level patterns recur across this set, although their magnitude varies by model.
GPT-5 \citep{singh2026openaigpt5card} and Claude~Opus~4.6 \citep{anthropic2026claudeopus46} remain to be evaluated. The same evaluation framework can accommodate these and future models.

\subsection{Cross-Encoders and Late-Interaction Retrieval}

Cross-encoder rerankers (models that jointly encode query and document to produce a relevance score) and late-interaction models such as ColBERT \citep{khattab2020colbert} occupy a middle ground between bi-encoder speed and LLM reasoning depth.
We evaluate this middle ground directly (\S\ref{sec:rerank}) by crossing four first-stage retrievers with cross-encoder and LLM listwise rerankers on BEIR and BRIGHT \citep{sun2023rankgpt,weller2025rank1testtimecomputereranking}. Reranking improves reasoning-heavy retrieval, while the strongest embedding alone remains best on semantic retrieval.

\citet{li2024ragvslcllms} compare RAG with long-context LLMs on QA tasks. Long-context LLMs achieve higher quality at greater token cost, while a Self-Route hybrid retains 93--100\% of their quality at 38--61\% of the token cost, depending on the model. This result supports the hybrid recommendation in \S\ref{sec:discussion}.


% ============================================================
\section{Extended Limitations}
\label{app:extended_limitations}
% ============================================================

We expand on the six limitations summarised in \S\ref{sec:limitations}.

\paragraph{(i) LLM coverage.}
The ten evaluated LLMs span six families, dense and mixture-of-experts architectures, and reasoning and instruct variants. They provide a snapshot of a fast-moving frontier.
At current public rates, the Pareto frontier includes Gemini~3.1~Pro and the leading embeddings; the lowest-cost LLM evaluated scores below the best embedding.
Closed APIs not yet covered (GPT-5 \citep{singh2026openaigpt5card} and Claude~Opus~4.6 \citep{anthropic2026claudeopus46}) may have different cost--capability profiles, so the numerical conclusions should be updated as new evaluations become available.
The released evaluation pipeline makes this straightforward to extend as new models are released.

\paragraph{(ii) Supervision asymmetry.}
Classification uses labelled-reference kNN for embeddings and zero-shot prompting for LLMs.
This setup reflects deployments with scarce labelled data and gives the two pipelines different supervision. Five examples preserve performance on the small-label tasks and reduce Banking77 performance (\S\ref{sec:fewshot}).

\paragraph{(iii) Small-corpus retrieval.}
The corpus-in-context protocol places entire corpora (82--415 documents) in the LLM prompt, a setting that grants the LLM three structural advantages: the corpus fits entirely in context, prompt caching amortises input cost across queries, and the model can reason globally over the full document set.
Each is a property of the protocol rather than a measured effect; we did not vary corpus size, so we do not quantify how retrieval quality would change at larger scale.
At BEIR or production scale, placing the full corpus in context is generally impractical; an indexed first stage is needed, so the small-corpus cost gap is a lower bound for this protocol.
We complement corpus-in-context retrieval with bi-encoder, cross-encoder, and LLM-reranker comparisons on BEIR and BRIGHT (\S\ref{sec:rerank}). Reranking benefits reasoning-heavy retrieval; the strongest embedding retains the lead on semantic retrieval \citep{sun2023rankgpt,weller2025rank1testtimecomputereranking}.
nDCG@$k$ results for the small-corpus setting are available in the released data files and preserve the paradigm-level ordering.

\paragraph{(iv) Reasoning-control granularity.}
Reasoning controls differ by provider: Gemini's \texttt{reasoning\_effort=low} provides a soft hint, while the open models support fully disabled reasoning.
We use both settings. Reduced reasoning preserves or improves retrieval for four of six models across five families (Figure~\ref{fig:cost_breakdown}b); token accounting varies by provider.

\paragraph{(v) Throughput comparison.}
We serve two open-weight LLMs (Qwen3.6-27B and Qwen3.6-35B-A3B) and the embedding models on the same H100, removing API rate limits and hardware choice as confounds.
The resulting 2.5--736$\times$ gap is consistent with the difference between autoregressive decoding and a single encoder pass.

\paragraph{(vi) Task weighting sensitivity.}
Our aggregate metric weights all \mteblm{} tasks equally regardless of dataset size, task difficulty, or practical importance.
Alternative weighting schemes (by dataset size, normalised z-scores, or task-family weighting) may shift aggregate conclusions.
Aggregate numbers should therefore be interpreted alongside the reported category- and task-level results.
Alternative aggregation methods such as the Borda count used by MMTEB \citep{enevoldsen2025mmtebmassivemultilingualtext} are applicable to \mteblm{} and may produce different paradigm orderings.


% ============================================================
\section{Cost and Throughput Details}
\label{app:tokens}
% ============================================================

This section reports token usage, embedding throughput, cost sensitivity, and the cost methodology underlying \S\ref{sec:cost_results}--\ref{sec:throughput_results}.

\subsection{Detailed LLM Token Usage}
\label{app:tokens_table}

Table~\ref{tab:llm_tokens} provides the complete token-level breakdown for each LLM across all \mteblm{} tasks.
For each model, we report input, standard-output, cached, reasoning, and total tokens, together with the cost decomposition. Gemini reports reasoning separately, whereas OpenRouter includes it in the completion count.
The reasoning-heavy LLMs generate 8--26M reasoning tokens (Qwen3.6-27B 26M, GLM-4.7 25M, Kimi-K2.6 22M, Gemini~3~Flash 14M, Pro 9M), compared with 1--3M standard-output tokens. Reasoning accounts for 28--81\% of total inference cost across these models: 81\% for Qwen3.6-27B and 28\% for DeepSeek-V4-Flash.
The non-reasoning Flash-Lite costs \$6.85 and ranks below most embeddings.

\input{tables/llm_tokens}

\subsection{Embedding Throughput and Cost}
\label{app:emb_throughput}

Table~\ref{tab:embedding_throughput} reports the throughput (tokens/second) and derived cost per million tokens for each of the 26 embedding models, benchmarked on a single NVIDIA H100 80GB HBM3 GPU at sequence length 512 with maximum batch size.
Costs range from \$0.00016/MTok (mE5-small, 118M parameters, 4.3M tok/s) to \$0.047/MTok (F2LLM-v2-14B, 14.7K tok/s).
Total benchmark cost spans 291$\times$ across the embedding models, compared with a 1{,}431$\times$ gap between the best embedding and Gemini~3.1~Pro.

Each model's total benchmark cost uses its own tokenizer's token count (4.4--5.5M tokens across vocabularies), avoiding tokenizer mismatch.

\input{tables/embedding_throughput}

\subsection{Cost Sensitivity Analysis}
\label{app:cost_sensitivity}

Table~\ref{tab:cost_sensitivity} shows how the LLM-to-embedding cost ratio varies under alternative hardware and pricing assumptions.
We consider five scenarios for embedding inference:
\begin{enumerate}
    \item \textbf{H100 spot} (\$2.49/hr, Lambda Labs, March 2026): baseline, 1{,}431$\times$
    \item \textbf{H100 on-demand} (\$3.99/hr): higher price, 893$\times$
    \item \textbf{A100 spot} (\$1.49/hr, est.\ 1.5$\times$ slower): 1{,}594$\times$
    \item \textbf{L4 spot} (\$0.49/hr, est.\ 3$\times$ slower): much lower throughput, 2{,}424$\times$
    \item \textbf{Commercial API} (e.g.\ \$0.10/MTok): fixed per-token pricing, 338$\times$
\end{enumerate}
In all five scenarios, the cost ratio exceeds 300$\times$.

\input{tables/cost_sensitivity}

\paragraph{Per-query economics and training amortisation.}
The benchmark-level ratios translate directly to per-query terms: both paradigms answer the same queries, so dividing by the query count rescales both sides equally and leaves the 29--1{,}431$\times$ range unchanged.
Two effects widen the gap further under sustained deployment.
First, embedding training is a one-time cost amortised over the deployment lifetime (and, for public checkpoints, across all users), so its per-query contribution vanishes at scale, whereas LLM inference cost recurs in full on every query.
Second, document embeddings are computed once and reused across queries, while corpus-in-context reading repeats per query and is only partially offset by prompt caching.
The benchmark-pass comparison is therefore conservative: it charges embeddings their full encoding cost while granting the LLM cached-input rates.

\subsection{Detailed Cost Methodology}
\label{app:cost_detail}

\paragraph{Embedding cost pipeline.}
For each model, we tokenize all \mteblm{} tasks with its HuggingFace \texttt{AutoTokenizer}. We then benchmark maximum-throughput processing on an NVIDIA H100 80GB at sequence length 512, using the largest batch that fits in memory. Cost is $(\text{tokens} / 10^6) \times (r_{\text{GPU}} / T)$, where $r_{\text{GPU}} = \$2.49$/hr (Lambda Labs spot pricing, March 2026) and $T$ is sustained throughput in tokens/hour.

Token counts range from 4.4M (Gemma-based tokenizers) to 5.5M (Mistral-based), a 26\% range driven by vocabulary differences.
Using a single proxy tokenizer (e.g., GPT-2 at 6.75M tokens) would overestimate embedding costs by 22--54\%.

\paragraph{LLM cost pipeline.}
Token usage is extracted from each provider's \texttt{usage\_stats} in every task's JSON result file.
For Gemini, \texttt{input\_tokens}\,=\,\texttt{prompt\_token\_count} (including cached tokens), \texttt{output\_tokens}\,=\,\texttt{candidates\_token\_count} (excluding reasoning), and \texttt{cached\_tokens}\,=\,\texttt{cached\_content\_token\_count}. Total tokens equal input plus output plus reasoning.
For the open models served via OpenRouter, the completion count already includes reasoning and \texttt{total}\,=\,input\,+\,output, with reasoning reported separately.
We bill every generated token, $\text{total}-\text{input}$, at the output rate. This equals output plus reasoning for Gemini and the reasoning-inclusive completion count for open models. Each model's reasoning share comes from its usage statistics (Figure~\ref{fig:cost_breakdown}).
For multilingual tasks (STS17, STS22v2, RTE3), per-language entries report different token counts and are summed; for classification tasks with multiple entries per split, all entries share identical \texttt{usage\_stats} and are counted once.

Gemini API prices per MTok input/output are \$0.25/\$1.50 for Flash-Lite, \$0.50/\$3.00 for Flash, and \$2.00/\$12.00 for Pro; cached input is billed at 10\%. Open models use OpenRouter public rates (March/June 2026).

\paragraph{Throughput measurement.}
We serve both paradigms on the same NVIDIA H100 and measure sustained tokens/second.
Open-weight LLMs (Qwen3.6-27B, Qwen3.6-35B-A3B) run under vLLM (BF16, tensor-parallel\,=\,1, 256 concurrent requests, 200/100 input/output tokens); embedding models run batch inference on the same GPU.
Holding hardware fixed removes API rate limits and GPU choice as explanations for the measured gap.


% ============================================================
\section{Reproducibility}
\label{app:reproducibility}
% ============================================================

All code, datasets, result files, and analysis scripts are released at \url{https://github.com/embeddings-benchmark/embedders-dilemma} and implemented with the MTEB framework.

\paragraph{LLM prompt templates.}
Prompt design follows a consistent structure across all five task categories:
\begin{itemize}
    \item \textbf{System prompt}: one sentence describing the task type and output format.
    \item \textbf{User prompt}: task-specific instruction, followed by the input text(s).
    \item \textbf{Output schema}: a Pydantic model with fields \texttt{reasoning} (free text, optional) and \texttt{output} (constrained to valid labels or numeric ranges). JSON mode is enabled for all calls.
\end{itemize}
Example system prompts: \emph{Classification}: ``You are a text classifier. Return a JSON with your reasoning and the predicted label from the provided list.'' \emph{Retrieval}: ``You are a retrieval assistant. Given a query and a set of numbered documents, identify the document IDs most relevant to the query. Return a JSON list of IDs in ranked order.''

\paragraph{Schema validation and failure handling.}
All LLM outputs are validated against the Pydantic schema after generation.
If validation fails (e.g., invalid label or malformed JSON), the prompt is retried with an appended correction instruction, up to twice per sample.
The final validation failure rate is $<$0.3\%; failed samples receive a null prediction and are counted as incorrect.

\paragraph{Token usage tracking.}
For each API call we log: \texttt{prompt\_token\_count}, \texttt{candidates\_token\_count}, \texttt{cached\_content\_token\_count}, \texttt{total\_token\_count}, and \texttt{thoughts\_token\_count} from \texttt{GenerateContentResponse.usage\_metadata}.
Thinking tokens are verified as $\texttt{total} - \texttt{prompt} - \texttt{candidates}$ to cross-check against \texttt{thoughts\_token\_count}; agreement was exact in all cases.
Pricing sources and throughput configurations are documented in \S\ref{app:cost_detail}.
